\section{Conclusion}
\label{sec:conclusion}

This development carries out a consistency argument for the encoded \BGA{}
proof system in \QGA{}, implemented as an extension of Isabelle/Pure.
The theorem \isa{BGA\_syntactically\_consistent} states that no two proof
codes are accepted, from empty hypotheses, for an equation and its
corresponding disequality.  The argument proves the checker sound and uses
determinism of evaluation to rule out such a pair.

Two features make the proof manageable.  The abstract encoding interface
separates the soundness argument from the arithmetic of a particular
G\"odel numbering; the Cantor-pairing construction then supplies that
interface.  Evaluation at a fixed fuel budget makes the satisfaction
premises used in the proof-list induction decidable, allowing the argument
to meet \QGA's habeas quid requirements.

An attempt to
extend the construction to \QGA{} itself is recorded in
Appendix~\ref{sec:self} for future collaborators.
